% !TeX root = RJwrapper.tex
\title{grip: Multiscale Graph Drawing and Layout Evaluation in R}


\author{by Pawel Gajer and Jacques Ravel}

\maketitle

\abstract{%
The \textbf{grip} package provides a comprehensive environment for
multiscale graph drawing in R, supporting scalable 2D and 3D layouts for
graphs ranging from tens to tens of thousands of vertices. The package
has evolved from an R implementation of the classical GRIP algorithm into
a family of multiscale layout engines spanning two distinct regimes:
combinatorial GRIP for unweighted graph structure and weighted GRIP for
graphs whose edge lengths encode geometry. The package offers a
consistent R API backed by a C++ core via Rcpp, validated parameter
presets for common graph families, a principled layout quality scoring
framework, reproducible comparison workflows, per-level trace
diagnostics, and a library of synthetic graph family generators covering
grids, manifolds, fractal structures, trees, and porous 3D bodies. This
article describes the package architecture, the user-facing APIs, the
diagnostic and benchmarking infrastructure, and demonstrates
representative workflows on synthetic and real-data examples.
}

\section{Introduction}\label{sec:introduction}

Visualizing graphs (networks) is a fundamental task in data analysis
across disciplines---from social network analysis to systems biology to
knowledge graph exploration. Force-directed layout algorithms, which
model vertices as charged particles and edges as springs, are the most
widely used family of methods for producing readable graph drawings.
Classical algorithms such as Fruchterman--Reingold
\citep{fruchterman1991} and Kamada--Kawai \citep{kamada1989} produce good
layouts for small graphs but scale poorly: their $O(n^2)$ or $O(n^3)$
per-iteration cost becomes prohibitive beyond a few hundred vertices.

For the large graphs routinely encountered in biological networks, social
media analysis, and knowledge graph construction, multiscale methods are
needed. The GRIP algorithm \citep{gajer2002, gajer2004} addresses this by
building a hierarchy of progressively coarser vertex subsets using
maximal independent set (MIS) filtrations, placing vertices
coarse-to-fine, and refining with local force-directed iterations at each
level. This yields near-linear time complexity in practice and produces
high-quality layouts even for graphs with thousands of vertices.

The R ecosystem offers several tools for graph layout. The \textbf{igraph}
package \citep{csardi2006} provides Fruchterman--Reingold, Kamada--Kawai,
and DrL implementations. The \textbf{ggraph} package \citep{pedersen2024}
adds a grammar-of-graphics layer but delegates layout computation to
igraph. The \textbf{graphlayouts} package \citep{schoch2024} implements
stress-based layouts with pivot MDS initialization. However, a gap
remains between layout \emph{generation} and layout \emph{assessment}: no
existing R package combines scalable multiscale layout computation with a
principled framework for quality evaluation, reproducible comparison
workflows, and diagnostic tracing.

The \textbf{grip} package addresses this gap. It began as an R
implementation of classical GRIP and has evolved into a broader
graph-drawing platform supporting:
%
\begin{itemize}
\item a family of multiscale layout engines for both unweighted
  (combinatorial) and weighted (geometry-aware) graphs,
\item validated parameter presets for common graph families,
\item a quality scoring framework that can evaluate layouts from any source,
\item trace and diagnostic tooling for understanding multiscale behavior,
\item a library of synthetic graph family generators for benchmarking, and
\item 2D and 3D layout support with interactive and static visualization.
\end{itemize}

The author of this package (P.\ Gajer) is one of the original developers
of the GRIP algorithm. This R implementation was motivated by the need
for a well-tested, accessible version of GRIP within the R ecosystem,
with the workflow tooling for quality assessment and reproducible layout
selection that was not part of the original publications.

The remainder of this article is organized as follows.
Section~\ref{sec:evolution} traces the evolution of the package from
a single-algorithm wrapper to a family of layout engines.
Section~\ref{sec:architecture} describes the package architecture.
Section~\ref{sec:apis} presents the user-facing layout APIs.
Section~\ref{sec:families} covers the graph family generators and
benchmark utilities. Section~\ref{sec:diagnostics} describes the
layout diagnostics and comparison framework.
Section~\ref{sec:examples} provides worked examples on synthetic and
real data. Section~\ref{sec:comparison} compares \textbf{grip} with
related R tools. Section~\ref{sec:discussion} discusses when to use
each layout variant, and Section~\ref{sec:conclusion} concludes.


\section{Evolution of the package}\label{sec:evolution}

The \textbf{grip} package has passed through three distinct stages of
development, each extending the scope of the layout engine while
maintaining backward compatibility with earlier APIs.

\paragraph{Stage 1: Classical GRIP.}
The initial release implemented the GRIP algorithm as described in
\citet{gajer2002} and \citet{gajer2004}: a maximal independent set (MIS)
filtration hierarchy, coarse-to-fine vertex placement, and local
force-directed refinement at each level. This engine is preserved as
\texttt{grip.layout.legacy()} for backward compatibility.

\paragraph{Stage 2: Quality-first GRIP with global repulsion.}
The current default engine, accessible as \texttt{grip.layout()} (and
its alias \texttt{grip.layout.globalrep()}), extends classical GRIP with
several improvements aimed at layout quality:
%
\begin{itemize}
\item coarse active-set global repulsion that reduces multiscale
  foldovers on structured graph families,
\item improved insertion-anchor strategies for vertex placement at each
  level, including distance-band and spread-based selection,
\item alternative final-stage refinement modes (Fruchterman--Reingold
  or KK-with-repulsion),
\item improved level-0 placement via optional local KK micro-polish, and
\item optional post-layout and in-core landmark geodesic KK (LGKK)
  refinement hooks.
\end{itemize}

\paragraph{Stage 3: Weighted geometry-aware GRIP.}
The most recent addition is a weighted sister family of GRIP methods
accessible through \texttt{grip.layout.weighted()} and its trace
counterpart \texttt{grip.layout.trace.weighted()}. Weighted GRIP
addresses a fundamental limitation of classical GRIP: because the
original algorithm constructs its MIS hierarchy using combinatorial
graph distance (hop count), it cannot preserve nontrivial edge-length
geometry. The weighted variants replace combinatorial operations with
geodesic-distance equivalents throughout the pipeline:
%
\begin{itemize}
\item weighted MIS filtration via \texttt{grip.build.misf.weighted()},
\item weighted neighborhood caches based on Dijkstra distances,
\item weighted insertion anchors using geodesic least-squares fitting, and
\item weighted local refinement forces proportional to edge-length targets.
\end{itemize}

The weighted API preserves the same parameter structure as the
unweighted API and adds its own family of validated presets (mesh,
cylinder, torus, sphere, irregular manifolds, tree, carpet). This
sister-API design ensures that the combinatorial engine remains
unchanged for users who do not need weighted layouts, while providing a
parallel path for geometry-aware graph drawing.

The package now supports two conceptually different multiscale layout
regimes: combinatorial GRIP for unweighted graph structure, and weighted
GRIP for graphs whose edge lengths encode geometry. This distinction is
central to the package design and is reflected in the API, the presets,
and the diagnostic tools.

\begin{figure}[ht]
\centering
\includegraphics[width=\linewidth]{../figures/misf-hierarchy}
\caption{MIS filtration hierarchy on a $6 \times 6$ mesh. Left: all 36
  vertices at the finest level. Center: an intermediate filtration level
  retaining a subset of vertices. Right: the coarsest level with only
  the initial seed vertices. Active vertices are shown as solid dots;
  inactive vertices as faint dots.}
\label{fig:misf}
\end{figure}


\section{Package architecture}\label{sec:architecture}

The \textbf{grip} package is organized in three layers: an R front-end
providing the public API, a C++ back-end implementing the
performance-critical multiscale engines, and a workflow layer providing
scoring, comparison, and diagnostic tools.

\subsection{R front-end}

The R layer provides idiomatic entry points to the layout engines,
parameter validation, preset resolution, disconnected-component
handling, and result packaging. All layout functions accept graphs in
either edge-list format (a two-column integer matrix) or adjacency-list
format (a list of integer neighbor vectors), with optional parallel
weight lists for weighted variants. The \texttt{seed} argument ensures
full reproducibility of layouts.

\subsection{C++ back-end}

The computational core is implemented in C++ and accessed via Rcpp
\citep{eddelbuettel2011}. This includes the MIS filtration construction,
the coarse-to-fine placement loop, the local force-directed refinement
iterations with adaptive temperature control, the coarse-level global
repulsion computation, and the weighted variants of all these operations.
The C++ layer also implements shortest-path computation (BFS for
unweighted, Dijkstra for weighted) and the optional LGKK refinement
engine.

\subsection{Workflow layer}

Beyond layout computation, the package provides a workflow layer for
systematic layout evaluation:
%
\begin{itemize}
\item \texttt{grip.score.layout()} computes quality metrics for
  individual layouts,
\item \texttt{grip.compare.layouts()} runs candidates across multiple
  seeds and aggregates scores,
\item \texttt{grip.layout.trace()} and \texttt{grip.layout.trace.weighted()}
  capture per-level snapshots with optional diagnostics, and
\item the graph family generators provide standardized test inputs for
  benchmarking.
\end{itemize}

\begin{figure}[ht]
\centering
\includegraphics[width=0.85\linewidth]{../figures/architecture}
\caption{Three-layer architecture of the \textbf{grip} package: R
  front-end (public API, presets, result packaging), C++ back-end
  (MIS filtration, force-directed engines, shortest paths), and
  workflow layer (scoring, comparison, diagnostics, graph generators).
  The workflow layer can score layouts from any source.}
\label{fig:architecture}
\end{figure}


\section{User-facing layout APIs}\label{sec:apis}

Table~\ref{tab:api-summary} summarizes the public layout functions
organized by task. All functions return an $n \times d$ numeric matrix
of vertex coordinates, where $d \in \{2, 3\}$.

\begin{table}[ht]
\centering
\caption{Public layout functions in the \textbf{grip} package.}
\label{tab:api-summary}
\begin{tabular}{lllc}
\toprule
Function & Engine & Graph type & Trace \\
\midrule
\texttt{grip.layout()} & Current quality-first & Unweighted & No \\
\texttt{grip.layout.trace()} & Current quality-first & Unweighted & Yes \\
\texttt{grip.layout.weighted()} & Weighted geometry-aware & Weighted & No \\
\texttt{grip.layout.trace.weighted()} & Weighted geometry-aware & Weighted & Yes \\
\texttt{grip.layout.legacy()} & Classical GRIP & Unweighted & No \\
\texttt{grip.layout.trace.legacy()} & Classical GRIP & Unweighted & Yes \\
\bottomrule
\end{tabular}
\end{table}

\subsection{Core layout interface}

The primary function \texttt{grip.layout()} accepts either an edge-list
matrix or an adjacency list:

\begin{verbatim}
library(grip)

# Edge-list input: 8x8 mesh in 2D
edges <- edges.mesh(8, 8)
coords <- grip.layout(edges, n = 64, dim = 2, preset = "mesh", seed = 1)

# Adjacency-list input with weights
adj <- list(c(2L), c(1L, 3L), c(2L, 4L), c(3L))
wts <- list(c(1.0), c(1.0, 2.0), c(2.0, 1.5), c(1.5))
coords <- grip.layout.weighted(adj_list = adj, weight_list = wts,
                                n = 4, dim = 2, seed = 12)
\end{verbatim}

Key parameters shared across all engines include \texttt{dim} (layout
dimension, 2 or 3), \texttt{placement} (initial arrangement:
\texttt{"barycenter"} or \texttt{"circle"}), \texttt{rounds} and
\texttt{final\_rounds} (refinement iteration counts), \texttt{num\_init}
(vertices at the coarsest level), \texttt{num\_nbrs} (local
neighborhood size), \texttt{r} and \texttt{s} (temperature adaptation
parameters), \texttt{repulsion\_factor} (finest-level repulsion
strength), and \texttt{seed} (for reproducibility).

Disconnected graphs are handled automatically by default
(\texttt{disconnected = "components"}): each connected component is laid
out independently and the results are packed into a single coordinate
matrix.

\subsection{Parameter presets}

Presets are validated parameter bundles for common graph families. Each
preset sets sensible defaults; any argument supplied explicitly by the
user overrides the corresponding preset value. The combinatorial engine
ships four presets (mesh, carpet, tree, torus), while the weighted engine
adds three more (cylinder, sphere, irregular manifolds). Table~\ref{tab:presets}
lists the available presets.

\begin{table}[ht]
\centering
\caption{Available parameter presets by engine.}
\label{tab:presets}
\begin{tabular}{llll}
\toprule
Preset & Engine & Target family & Tuned on \\
\midrule
\texttt{"mesh"} & Both & Rectangular grid & $8 \times 8$ to $12 \times 12$ \\
\texttt{"carpet"} & Both & Sierpinski carpet & Level 3--4 \\
\texttt{"tree"} & Both & Tree-like graphs & Binary trees, depth 5--6 \\
\texttt{"torus"} & Both & 3D torus/cylinder & $8 \times 8$ to $20 \times 20$ \\
\texttt{"cylinder"} & Weighted & Cylinder meshes & Various sizes \\
\texttt{"sphere"} & Weighted & Sphere meshes & Triangulated spheres \\
\texttt{"irregular"} & Weighted & Irregular manifolds & Irregular surfaces \\
\bottomrule
\end{tabular}
\end{table}

\subsection{Multiscale trace}

The trace functions capture snapshots of the layout at intermediate
stages of the multiscale process, supporting animated visualizations
and diagnostic inspection. The \texttt{trace} parameter controls
granularity (\texttt{"level"} for coarser or \texttt{"round"} for
finest snapshots), and optional per-frame diagnostics include
edge-length CV, non-edge separation, stress, and Procrustes RMSE
against a target layout.

\begin{verbatim}
edges <- edges.mesh(6, 6)
tr <- grip.layout.trace(edges, n = 36, dim = 2,
                         preset = "mesh", seed = 31,
                         trace = "level")
\end{verbatim}

Each frame in \texttt{tr\$frames} contains coordinates for all
vertices, with \texttt{NA} entries for vertices not yet introduced at
that level. Figure~\ref{fig:trace} shows four stages of a mesh layout,
from the initial coarse placement through progressive refinement.

\begin{figure}[ht]
\centering
\includegraphics[width=\linewidth]{../figures/trace-mesh-6x6}
\caption{Multiscale trace of a $6 \times 6$ mesh layout. From left to
  right: coarsest initialization, early refinement, late refinement,
  and final layout. Vertices not yet introduced appear as faint dots.}
\label{fig:trace}
\end{figure}

\subsection{LGKK refinement hooks}

Both the combinatorial and weighted engines support optional
landmark geodesic KK (LGKK) refinement. LGKK uses a sparse set of
farthest-point landmarks per vertex to compute approximate geodesic
distances, then applies Kamada--Kawai-style gradient steps to improve
geodesic fidelity. The integration is available as:
%
\begin{itemize}
\item \texttt{lgkk\_polish\_rounds}: post-layout polish after the
  multiscale solve completes,
\item \texttt{lgkk\_rounds\_coarse}, \texttt{lgkk\_rounds\_pre\_final},
  \texttt{lgkk\_rounds\_final}: in-core LGKK refinement at different
  stages of the multiscale hierarchy.
\end{itemize}

The LGKK subsystem is also available as standalone functions
(\texttt{grip.prepare.landmark.geodesic.kk()},
\texttt{grip.score.landmark.geodesic.kk()},
\texttt{grip.optimize.landmark.geodesic.kk()}) for users who wish to
apply geodesic refinement to layouts from any source.

\subsection{Plotting}

The \texttt{grip.plot()} function provides a lightweight plotting helper.
For 2D layouts, it draws vertices and edges using base graphics. For 3D
layouts, it supports interactive viewing via the optional \textbf{rgl}
package (\texttt{projection = "rgl"}) and static orthographic projection
(\texttt{projection = "ortho"}) suitable for publications. The
\texttt{grip.project.3d()} function applies an orthographic rotation
defined by azimuth and elevation angles, returning a 2D projection
matrix.


\section{Graph families, geometry, and benchmark utilities}\label{sec:families}

A distinctive feature of \textbf{grip} is its library of synthetic graph
family generators, which serves both as a testing substrate and as a
benchmarking resource. The package provides 29 edge-generator functions
organized into several categories (Table~\ref{tab:families}).

\begin{table}[ht]
\centering
\caption{Graph family generators in the \textbf{grip} package. Each
function returns a two-column integer edge matrix. Families marked
with~$\dagger$ also provide companion \texttt{.surface.embedding()}
functions returning vertex coordinates.}
\label{tab:families}
\begin{tabular}{ll}
\toprule
Category & Generators \\
\midrule
Basic & \texttt{edges.path()}, \texttt{edges.cycle()} \\
Regular grids & \texttt{edges.mesh()}, \texttt{edges.cylinder()}, \\
              & \texttt{edges.torus()}, \texttt{edges.sphere()} \\
Trees & \texttt{edges.kary.tree()} \\
Irregular manifolds$^\dagger$ & \texttt{edges.irregular.ball()}, \texttt{edges.irregular.torus()}, \\
                               & \texttt{edges.irregular.sphere()}, \texttt{edges.irregular.annulus()}, \\
                               & \texttt{edges.irregular.pair.of.pants()}, \\
                               & \texttt{edges.irregular.double.torus()}, \\
                               & \texttt{edges.irregular.shell()} \\
Fractal/recursive & \texttt{edges.sierpinski.triangle()}, \\
                  & \texttt{edges.sierpinski.tetrahedron()}, \\
                  & \texttt{edges.sierpinski.carpet()}, \\
                  & \texttt{edges.vicsek()}, \texttt{edges.menger.sponge()} \\
Porous 3D bodies$^\dagger$ & \texttt{edges.cube.periodic.tunnels()}, \\
                            & \texttt{edges.cube.asymmetric.cavities()}, \\
                            & \texttt{edges.cube.channel.network()} \\
Polyhedra & \texttt{edges.cube()}, \texttt{edges.triangulated.polyhedron()}, \\
          & \texttt{edges.triangulated.annulus()}, \\
          & \texttt{edges.triangulated.pair.of.pants()} \\
\bottomrule
\end{tabular}
\end{table}

The irregular manifold and porous 3D body families are particularly
important for benchmarking weighted GRIP, because they provide graphs
with nontrivial edge-length geometry and known ground-truth embeddings
(via the companion \texttt{.surface.embedding()} functions). This
enables quantitative evaluation of how well a layout algorithm preserves
the intrinsic geometry of the underlying surface.

\begin{figure}[ht]
\centering
\includegraphics[width=\linewidth]{../figures/preset-comparison}
\caption{Effect of validated parameter presets on a level-3 Sierpinski
  carpet. Left: default parameters. Right: the \texttt{"carpet"} preset,
  which produces substantially lower stress and more uniform edge lengths.}
\label{fig:preset}
\end{figure}

The fractal families (Sierpinski triangle, carpet, tetrahedron; Vicsek;
Menger sponge) and the recursive mask generators
(\texttt{edges.recursive.mask.grid()},
\texttt{edges.recursive.triangle.mask()},
\texttt{edges.recursive.cube.mask()}) provide graphs with challenging
self-similar structure at controlled scales, useful for stress-testing
multiscale algorithms.


\section{Layout diagnostics and comparison}\label{sec:diagnostics}

For real-world graphs---social networks, biological networks, knowledge
graphs---there is typically no ground-truth embedding. Layout quality
must be assessed through proxy metrics. The \textbf{grip} package
provides a systematic workflow for this task.

\subsection{Quality metrics}

The \texttt{grip.score.layout()} function computes six metrics for any
layout, regardless of how it was produced:

\begin{table}[ht]
\centering
\caption{Layout quality metrics available in \texttt{grip.score.layout()}.}
\label{tab:metrics}
\begin{tabular}{lll}
\toprule
Metric & Measures & Better \\
\midrule
Sampled stress & Fidelity to graph distances & Lower \\
Edge-length CV & Uniformity of edge lengths & Lower \\
Non-edge separation & Spacing of non-neighbors & Higher \\
Edge crossings & Readability (2D only) & Lower \\
Cluster separation & Group visual separation & Higher \\
Procrustes stability & Seed-to-seed reproducibility & Lower \\
\bottomrule
\end{tabular}
\end{table}

Sampled stress measures the normalized RMSE between optimally scaled
Euclidean distances and graph-theoretic distances for a random sample of
vertex pairs (BFS for unweighted, Dijkstra for weighted). Edge-length
CV reports the coefficient of variation of all edge lengths. Non-edge
separation computes the ratio of the minimum sampled non-edge distance
to the median edge length. Cluster separation (optional, requires
external labels) reports the ratio of mean inter-cluster to mean
intra-cluster distances. Procrustes stability aligns layouts from
different seeds via Procrustes analysis and reports mean RMSE.

\subsection{Comparison workflow}

The \texttt{grip.compare.layouts()} function provides a systematic
comparison pipeline:

\begin{verbatim}
edges <- edges.mesh(6, 6)
cmp <- grip.compare.layouts(
  edges = edges, n = 36, dim = 2,
  candidates = c("default", "mesh"),
  seeds = 1:3
)
\end{verbatim}

This runs each candidate configuration across multiple seeds, computes
per-run scores, and aggregates them into a summary table with
rank-based composite scores. The function also supports a
\texttt{search} interface for local parameter grid searches around a
baseline configuration, enabling systematic exploration of the
parameter space.

The composite score uses configurable weights (default: stress 0.35,
edge-length CV 0.20, non-edge separation 0.20, Procrustes stability
0.15, edge crossings 0.05, cluster separation 0.05). It is useful as a
screening tool for shortlisting candidates but should not be treated as
the sole basis for layout selection.

\subsection{Reproducibility}

The \texttt{seed} parameter in all layout functions ensures that
layouts are fully reproducible. The comparison workflow records seeds,
parameter settings, and all metric values, enabling scripted pipelines
that produce identical results across runs. The trace functions
additionally record intermediate states, supporting detailed
post-mortem analysis of the multiscale process.


\section{Worked examples}\label{sec:examples}

\subsection{Synthetic example: combinatorial versus weighted GRIP on a
torus}\label{sec:ex-torus}

The torus is a natural test case for weighted GRIP, because its
intrinsic geometry assigns meaningful lengths to edges that the
combinatorial engine ignores. We compare the two engines on a weighted
$12 \times 12$ torus in 3D:

\begin{verbatim}
edges <- edges.torus(12, 12)
n <- max(edges)

# Combinatorial GRIP
coords.comb <- grip.layout(edges, n = n, dim = 3,
                            preset = "torus", seed = 1)

# Weighted GRIP (using unit edge weights for comparison)
edge_weights <- rep(1.0, nrow(edges))
coords.wt <- grip.layout.weighted(edges, n = n, dim = 3,
                                   edge_weights = edge_weights,
                                   preset = "torus", seed = 1)
\end{verbatim}

For graphs with nontrivial edge-length geometry (e.g., meshes on curved
surfaces), the weighted engine preserves geodesic distances that the
combinatorial engine cannot represent. This distinction becomes
significant for irregular manifolds and porous 3D bodies where
edge lengths vary substantially.

\begin{figure}[ht]
\centering
\includegraphics[width=\linewidth]{../figures/triptych-saddle}
\caption{Combinatorial versus weighted GRIP on a $5 \times 5$ saddle
  mesh in 3D. Left: the target surface geometry. Center: combinatorial
  GRIP, which ignores edge-length information. Right: weighted GRIP,
  which preserves the intrinsic geodesic distances.}
\label{fig:triptych}
\end{figure}

\subsection{Real-data example: HMP microbial network}\label{sec:ex-hmp}

The package bundles a coarsened real-world graph derived from the Human
Microbiome Project (HMP) 16S amplicon analysis \citep{hmp2012}. The
original giant connected component contained 6,474 vertices from an
iterative $k$-nearest-neighbor graph with $k=3$ built on a $\geq 1\%$
relative-abundance OTU table. Greedy graph coarsening in two rounds
reduced this to 1,828 supernodes, preserving the large-scale topology
including arm-like extensions corresponding to distinct community state
types (CSTs).

This dataset is large enough that interactive manual layout selection is
impractical, making the systematic comparison workflow especially
valuable. The full pipeline is documented in the package's HMP vignette.

\begin{verbatim}
data(hmp.u01.gc.coarse)
cat("Vertices:", length(hmp.u01.gc.coarse$adj_list), "\n")
#> Vertices: 1828
cat("CST levels:", paste(sort(unique(
  hmp.u01.gc.coarse$vertex_data$cst)), collapse = ", "), "\n")
#> CST levels: I, II, III, IV-A, IV-B, IV-C, V
\end{verbatim}

\begin{figure}[ht]
\centering
\includegraphics[width=\linewidth]{../figures/hmp-layouts}
\caption{Coarsened HMP microbial network (1,828 vertices) in 3D,
  colored by community state type (CST). Three preset candidates are
  compared: default, tree, and a local-search candidate. The tree and
  local-search layouts show clearer CST separation in the arm-like
  extensions of the graph.}
\label{fig:hmp}
\end{figure}

\subsection{3D versus 2D: when the extra dimension matters}\label{sec:ex-3d}

For graphs with inherently three-dimensional structure---tori, spheres,
porous 3D bodies---forcing the layout into 2D necessarily introduces
distortion. The \textbf{grip} package supports 3D layout natively, and
the trace diagnostics can quantify the quality loss from dimensional
reduction.

\begin{verbatim}
edges <- edges.torus(8, 8)
n <- max(edges)

coords.2d <- grip.layout(edges, n = n, dim = 2,
                          preset = "torus", seed = 1)
coords.3d <- grip.layout(edges, n = n, dim = 3,
                          preset = "torus", seed = 1)

score.2d <- grip.score.layout(coords.2d, edges = edges, n = n)
score.3d <- grip.score.layout(coords.3d, edges = edges, n = n)
\end{verbatim}

In typical results, the 3D layout achieves substantially lower stress
and higher non-edge separation than the 2D layout on graphs with 3D
intrinsic structure, confirming that the extra dimension is not merely
cosmetic.

\begin{figure}[ht]
\centering
\includegraphics[width=\linewidth]{../figures/2d-vs-3d}
\caption{Weighted GRIP on a $5 \times 5$ saddle mesh in 2D (left) and
  3D (right). The 3D layout captures the intrinsic curvature of the
  surface, while the 2D layout necessarily introduces distortion.}
\label{fig:2d3d}
\end{figure}


\section{Comparison with related R tools}\label{sec:comparison}

Table~\ref{tab:pkg-comparison} provides a practical comparison of
\textbf{grip} with the three most relevant R packages for graph layout.

\begin{table}[ht]
\centering
\caption{Feature comparison of R packages for graph layout.}
\label{tab:pkg-comparison}
\begin{tabular}{lcccc}
\toprule
Feature & \textbf{grip} & \textbf{igraph} & \textbf{ggraph} & \textbf{graphlayouts} \\
\midrule
Multiscale layout & Yes & DrL only & No & No \\
3D layout & Yes & Limited & No & No \\
Weighted geometry-aware & Yes & No & No & No \\
Quality scoring & Yes & No & No & No \\
Comparison workflow & Yes & No & No & No \\
Trace diagnostics & Yes & No & No & No \\
Graph family generators & Yes & Some & No & No \\
Parameter presets & Yes & No & No & No \\
\bottomrule
\end{tabular}
\end{table}

The comparison should be understood as complementary rather than
competitive. The \textbf{igraph} package provides a broader range of
graph algorithms and analysis tools. The \textbf{ggraph} package excels
at publication-quality static visualization using the grammar of
graphics. The \textbf{graphlayouts} package provides fast stress-based
layouts. The \textbf{grip} package focuses specifically on the layout
computation and quality assessment problem, and its scoring functions
can evaluate layouts computed by any package.

To illustrate this complementarity, the quality metrics from
\texttt{grip.score.layout()} can assess layouts produced by
\textbf{igraph}'s algorithms:

\begin{verbatim}
# Score an igraph layout using grip metrics
edges <- edges.mesh(8, 8)
ig <- igraph::graph_from_edgelist(edges, directed = FALSE)
coords.kk <- igraph::layout_with_kk(ig)
score.kk <- grip.score.layout(coords.kk, edges = edges, n = 64)
\end{verbatim}

On small 2D social networks, classical igraph layouts (FR, KK) are
typically stronger on stress and edge crossings, consistent with their
design goals. The \textbf{grip} package's strengths emerge on larger
graphs, structured graph families, 3D layouts, and in workflows
requiring systematic parameter selection and quality assessment.

\begin{figure}[ht]
\centering
\includegraphics[width=\linewidth]{../figures/comparison-panel-mesh}
\caption{Layout comparison on a $12 \times 12$ mesh. Top row: FR, KK,
  and DrL from \textbf{igraph}. Bottom row: stress minimization from
  \textbf{graphlayouts}, and \textbf{grip} with default and mesh-preset
  parameters. The mesh preset produces near-perfect grid recovery with
  uniform edge lengths.}
\label{fig:comparison-mesh}
\end{figure}

\begin{figure}[ht]
\centering
\includegraphics[width=\linewidth]{../figures/comparison-panel-carpet}
\caption{Layout comparison on a level-4 Sierpinski carpet. Same method
  arrangement as Figure~\ref{fig:comparison-mesh}. The carpet preset
  preserves the recursive hole structure more faithfully than
  general-purpose methods.}
\label{fig:comparison-carpet}
\end{figure}


\section{Discussion}\label{sec:discussion}

The \textbf{grip} package now supports two distinct multiscale layout
regimes, and users should choose between them based on their data:

\paragraph{When to use combinatorial GRIP.}
The default \texttt{grip.layout()} is appropriate for unweighted graphs
or graphs where edge weights do not represent meaningful geometric
distances. This includes social networks, citation networks, and
biological interaction networks where edges represent presence or
absence of a relationship. The combinatorial engine is faster because
it uses BFS rather than Dijkstra for distance computation.

\paragraph{When to use weighted GRIP.}
The weighted engine \texttt{grip.layout.weighted()} is appropriate when
edge weights represent distances or similarities that should be
preserved in the layout. This includes mesh graphs from surface
discretizations, $k$-nearest-neighbor graphs built from distance
matrices, and any graph where the edge lengths encode geometry.

\paragraph{Sister APIs versus unified interface.}
The package deliberately maintains separate APIs for the combinatorial
and weighted engines rather than adding a mode switch to a single
function. This design ensures backward compatibility, makes the
distinction explicit in user code, and allows each engine to have its
own validated presets and default parameters.

\paragraph{Limitations.}
The current implementation does not exploit GPU parallelism, so very
large graphs (100,000+ vertices) may require substantial computation
time. The edge crossing count is exact but $O(m^2)$, limiting its use
to moderately sized graphs. The quality metrics are heuristic proxies,
not formal optimality guarantees.

\paragraph{Relation to companion papers.}
The algorithmic novelty of the weighted GRIP variants---the weighted
MISF construction, the geodesic refinement integration, and the
systematic benchmark results---is the subject of a separate methods
paper. The geodesic MDS functionality in the package
(\texttt{grip.prepare.geodesic.mds()} and related functions) addresses
a different use case and will be described separately.


\section{Conclusion}\label{sec:conclusion}

The \textbf{grip} package has evolved from an R implementation of the
GRIP algorithm into a comprehensive environment for multiscale graph
drawing. It now supports both combinatorial and geometry-aware
multiscale layouts in 2D and 3D, provides a principled quality
assessment framework, offers reproducible comparison workflows, and
includes a rich library of synthetic graph families for benchmarking.
The consistent API design, the C++ core for performance, and the
diagnostic tooling make \textbf{grip} a practical platform for both
applied graph visualization and graph-drawing research in R.

The package source code is available at
\url{https://github.com/pgajer/grip} and is being prepared for CRAN
submission.

\section{Acknowledgments}\label{sec:acknowledgments}

The GRIP algorithm was originally developed in collaboration with
Stephen G.\ Kobourov and Michael T.\ Goodrich. The HMP microbial
network data were derived from publicly available Human Microbiome
Project resources. Research reported in this publication was supported
by grant number INV-048956 from the Gates Foundation.

\bibliography{grip-paper}

\address{%
Pawel Gajer\\
Center for Advanced Microbiome Research and Innovation (CAMRI),
Institute for Genome Sciences, Department of Microbiology and
Immunology, University of Maryland School of Medicine\\%
Baltimore, Maryland, USA\\
%
\href{mailto:pgajer@gmail.com}{\nolinkurl{pgajer@gmail.com}}%
}

\address{%
Jacques Ravel\\
Center for Advanced Microbiome Research and Innovation (CAMRI),
Institute for Genome Sciences, Department of Microbiology and
Immunology, University of Maryland School of Medicine\\%
Baltimore, Maryland, USA\\
%
}
